\section{System Design}
\label{sec:design}

Pigs is structured as a three-layer system: a \emph{proxy} layer that routes HTTP requests, a \emph{phased runtime} layer that orchestrates the Pre$\rightarrow$Executor$\rightarrow$Post state machine, and a \emph{CLI} layer that provides an interactive agent interface. This section describes the core design of each layer.

\subsection{Architecture Overview}
\label{sec:overview}

The system operates as a single-process server listening on one port (default 3927). All three API protocols are served from the same port, distinguished by request path:

\begin{itemize}[nosep]
    \item \texttt{/chat/completions} $\rightarrow$ OpenAI Chat
    \item \texttt{/v1/messages} $\rightarrow$ Anthropic Messages
    \item \texttt{/responses} $\rightarrow$ OpenAI Responses
\end{itemize}

Requests are routed by model identifier: models without the \texttt{-pig} suffix are passed through to the upstream provider (with retry, body cleaning, and thinking-effort injection); models with the \texttt{-pig} suffix enter the phased runtime. Each configured model is automatically exposed as two identifiers: \texttt{\{model\}} (passthrough) and \texttt{\{model\}-pig} (phased).

\Cref{fig:arch} illustrates the request flow. When a \texttt{-pig} request arrives, the proxy constructs a \texttt{HttpRequestEnvelope} and hands it to the \texttt{HttpPhasedRuntime}. Each phase subrequest is sent back to the proxy via HTTP loopback (with an internal token header), re-entering the passthrough path with the \texttt{-pig} suffix removed. This design ensures that phase subrequests transparently inherit all proxy-level features.

\begin{figure}[t]
\centering
\begin{tikzpicture}[
  node distance=0.8cm,
  box/.style={draw, rounded corners, minimum width=3.5cm, minimum height=0.8cm, align=center, font=\small},
  client/.style={box, fill=blue!10},
  proxy/.style={box, fill=green!10},
  runtime/.style={box, fill=orange!10},
  provider/.style={box, fill=red!10},
  arr/.style={-Stealth, thick},
  loopback/.style={-Stealth, thick, dashed, draw=purple!70},
]
\node[client] (client) {Client (model: X-pig)};
\node[proxy, below=of client] (proxy1) {pigs-proxy\\(route by -pig)};
\node[runtime, below=of proxy1] (runtime) {HttpPhasedRuntime\\Pre $\to$ Executor $\to$ Post};
\node[proxy, below=of runtime] (proxy2) {pigs-proxy\\(passthrough, model: X)};
\node[provider, below=of proxy2] (provider) {Upstream Provider};

\draw[arr] (client) -- (proxy1);
\draw[arr] (proxy1) -- (runtime);
\draw[loopback] (runtime) -- node[right, font=\scriptsize, text=purple!70] {HTTP loopback} (proxy2);
\draw[arr] (proxy2) -- (provider);
\end{tikzpicture}
\caption{Request flow for a \texttt{-pig} model. Phase subrequests re-enter the proxy via HTTP loopback (dashed), inheriting passthrough features.}
\label{fig:arch}
\end{figure}
\subsection{Phase State Machine}
\label{sec:state}

The core of Pigs is a pure state machine with three phases and three transition signals. The state is encapsulated in \texttt{OrchestrationState}, which tracks the current phase, accumulated outputs, failure history, and iteration counters.

\begin{algorithm}[t]
\caption{Phase state machine transition}
\label{alg:state}
\begin{algorithmic}[1]
\Function{advance}{output}
  \State $phase \gets \text{current phase}$
  \If{$phase = \text{Pre}$}
    \If{$\text{detect}(output) = \texttt{PIGEND}$}
      \State \Return Complete \Comment{Short-circuit: trivially simple task}
    \ElsIf{$\text{detect}(output) = \texttt{PIGFAIL}$}
      \State \Return replan(output) \Comment{Failure: return to Pre}
    \Else
      \State $\text{pre\_output} \gets \text{strip}(output)$
      \State \Return Continue(Executor)
    \EndIf
  \ElsIf{$phase = \text{Executor}$}
    \State $\text{executor\_outputs} \gets \text{strip}(output)$
    \State \Return Continue(Post) \Comment{Always proceed to Post}
  \Else \Comment{$phase = \text{Post}$}
    \If{$\text{detect}(output) = \texttt{PIGEND}$}
      \State \Return Complete
    \ElsIf{$\text{detect}(output) = \texttt{PIGFAIL}$}
      \State \Return replan(output)
    \Else
      \If{$\text{post\_iterations} \geq \text{max}$}
        \State \Return Error(PostBudgetExceeded)
      \EndIf
      \State $\text{post\_iterations} \mathrel{+}= 1$
      \State \Return Continue(Post) \Comment{Stay in Post}
    \EndIf
  \EndIf
\EndFunction
\end{algorithmic}
\end{algorithm}

The state machine has several key properties:

\textbf{Markerless Post stays in Post.} When the Post phase produces output without a control marker, the state machine continues in Post (not back to Executor). This design choice prevents infinite Executor$\leftrightarrow$Post loops: the model is expected to make incremental progress within Post until it either completes or fails.

\textbf{Budget exhaustion is an error, not success.} When the Post iteration budget is exhausted, the state machine returns an \texttt{OrchestrationError}, not a successful completion. This prevents the system from silently accepting incomplete work.

\textbf{Failure preserves history.} When \texttt{PIGFAIL} triggers a replan, the failure output is preserved and presented to the next Pre phase as numbered failure records (``1st failure: \ldots'', ``2nd failure: \ldots''). This allows the model to learn from previous attempts.

\begin{figure}[t]
\centering
\begin{tikzpicture}[
  node distance=2cm,
  state/.style={draw, circle, minimum size=1.6cm, align=center, font=\small},
  terminal/.style={draw, rounded corners, minimum width=1.8cm, minimum height=0.8cm, align=center, font=\small, fill=gray!10},
  arr/.style={-Stealth, thick},
  every label/.style={font=\scriptsize},
]
\node[state, fill=blue!10] (pre) {Pre};
\node[state, fill=orange!10, right=of pre] (exec) {Exec.};
\node[state, fill=green!10, right=of exec] (post) {Post};
\node[terminal, right=of post] (done) {Complete};
\node[terminal, below=1.2cm of post, fill=red!10] (error) {Error};

\draw[arr] (pre) -- node[above] {no marker} (exec);
\draw[arr] (exec) -- node[above] {always} (post);
\draw[arr] (post) -- node[above] {\texttt{PIGEND}} (done);
\draw[arr] (pre) -- node[above, sloped] {\texttt{PIGEND}} (done);
\draw[arr, dashed] (post.south) -- node[right] {\texttt{PIGFAIL}} (pre.south);
\draw[arr, dashed] (post) -- node[right, sloped] {budget} (error);
\draw[arr, loop above] (post) to node[above] {no marker} (post);
\end{tikzpicture}
\caption{Phase state machine transitions. Solid arrows are forward transitions; dashed arrows are failure/budget paths. The self-loop on Post represents markerless continuation.}
\label{fig:state}
\end{figure}

\Cref{fig:state} illustrates the complete state machine. The Pre phase can short-circuit to Complete (trivially simple tasks), and the Post phase has a self-loop for markerless continuation.

\subsection{Control Markers}
\label{sec:markers}

Control markers are short text tokens that the model appends to its output to signal phase transitions:

\begin{itemize}[nosep]
    \item \texttt{PIGEND}: signals successful completion. The turn ends, and the accumulated visible text is returned to the client.
    \item \texttt{PIGFAIL}: signals that the current approach has failed and the task should be replanned from Pre.
\end{itemize}

Marker detection follows two rules:

\textbf{Last non-empty line only.} Only the final non-empty line of the model's output is checked for a marker. Earlier occurrences (e.g., the marker appearing in a sentence like ``the system outputs PIGEND when done'') are ignored. This prevents false positives from prose mentions.

\textbf{Reason required.} A marker is valid only if at least one non-marker line precedes it. A bare \texttt{PIGEND} with no preceding content is rejected, forcing the model to provide a reason before signaling completion.

Markers are stripped from the visible output via \texttt{strip\_markers}, which removes exact control lines while preserving prose that happens to contain marker words.

\subsection{Protocol-Native Request Envelope}
\label{sec:envelope}

The \texttt{HttpRequestEnvelope} is the central data structure that preserves the complete protocol-native request:

\begin{lstlisting}
pub struct HttpRequestEnvelope {
    pub method: String,           // HTTP method
    pub path_and_query: String,   // path + query string
    pub headers: Vec<HeaderPair>, // raw header bytes
    pub body: Value,              // complete JSON body
    pub protocol: Protocol,       // Chat / Anthropic / Responses
    pub client_model: String,     // e.g. "gpt-4-pig"
    pub real_model: String,       // e.g. "gpt-4" (-pig stripped)
    pub stream: bool,             // streaming requested?
    // ... internal fields for current user location,
    //     tool result IDs, tool result groups
}
\end{lstlisting}

Each phase clones the envelope and applies \emph{point mutations}:

\begin{enumerate}[nosep]
    \item \textbf{Model}: replaced with the real (non-\texttt{-pig}) model name.
    \item \textbf{Current user input}: appended with the phase-specific prompt suffix (Pre questions, Executor plan, Post review instructions).
    \item \textbf{Phase transcript}: native assistant/tool messages from prior phases are appended to the message history.
    \item \textbf{Stream}: set to \texttt{true} for streaming phase subrequests.
\end{enumerate}

All other fields---system/instructions, historical messages, tools/tool\_choice, reasoning parameters, cache controls, metadata, media blocks, and unknown provider extensions---remain untouched. This ensures that provider-specific features (e.g., Anthropic's \texttt{cache\_control}, OpenAI's \texttt{reasoning\_effort}) pass through without modification.

\subsection{HTTP Loopback Transport}
\label{sec:loopback}

Each phase subrequest is sent back to the local proxy via HTTP, not through an in-process function call. The \texttt{LoopbackPhaseTransport} implements the \texttt{PhaseTransport} trait:

\begin{itemize}[nosep]
    \item Converts wildcard listen addresses (\texttt{0.0.0.0}, \texttt{[::]}) to loopback (\texttt{127.0.0.1}, \texttt{::1}).
    \item Disables environment HTTP proxy settings (\texttt{.no\_proxy()}) to prevent external proxy interception.
    \item Adds a random internal token header (\texttt{x-pigs-phase-loopback}) that the proxy validates to identify legitimate phase subrequests.
    \item Strips hop-by-hop headers and the internal token before forwarding to the upstream provider.
\end{itemize}

This design has two benefits. First, phase subrequests transparently inherit all proxy-level features: channel selection, model mapping, body cleaning (e.g., removing empty content messages), thinking-effort injection, and retry logic. Second, the control flow is explicit and debuggable---each phase subrequest appears in proxy logs as a normal HTTP request, and the internal token distinguishes it from client requests.

\subsection{Bounded In-Memory Continuation}
\label{sec:continuation}

When the model returns tool calls during a phase, the runtime pauses and returns the tool calls to the client (or external agent). The complete phase state---including the orchestration state, phase transcript, visible text, and pending tool IDs---is stored in a \texttt{ContinuationStore}:

\begin{itemize}[nosep]
    \item \textbf{TTL}: each continuation expires after a configurable duration (default 30 minutes).
    \item \textbf{Capacity}: at most $N$ continuations are retained (default 256); the oldest is evicted when capacity is reached.
    \item \textbf{Tombstones}: when a continuation is expired, evicted, or consumed, a tombstone is placed for each tool ID, enabling distinct error reporting (Expired vs.\ Evicted vs.\ Consumed vs.\ Unknown).
    \item \textbf{No auth headers}: authentication headers are cleared before storing; the continuation uses the incoming request's headers on resume.
\end{itemize}

When the next request arrives with tool results, the runtime matches them by tool-call ID, validates protocol and model consistency, and resumes the paused phase. Partial parallel results (multiple tool calls returning in separate requests) are accumulated until all pending IDs are present.

\subsection{Continuous SSE Output}
\label{sec:sse}

For streaming requests, the runtime produces a continuous SSE stream that concatenates visible text from all phases. A \emph{marker line buffer} holds back the last non-empty line of each phase's output until the phase completes, then checks for control markers. If a marker is found, the held-back line is stripped; otherwise, it is forwarded. This prevents \texttt{PIGEND}/\texttt{PIGFAIL} from leaking into the client-visible stream, even when markers are split across network chunks.

The outer SSE encoder generates protocol-native event sequences for all three protocols:
\begin{itemize}[nosep]
    \item \textbf{Chat}: \texttt{chat.completion.chunk} deltas with role/content/tool\_calls, terminated by \texttt{finish\_reason} + \texttt{[DONE]}.
    \item \textbf{Anthropic}: \texttt{message\_start} $\rightarrow$ \texttt{content\_block\_start/delta/stop} $\rightarrow$ \texttt{message\_delta} $\rightarrow$ \texttt{message\_stop}.
    \item \textbf{Responses}: \texttt{response.created} $\rightarrow$ \texttt{output\_item/content\_part/text\_delta/done} $\rightarrow$ \texttt{response.completed}, with monotonic sequence numbers.
\end{itemize}

Error streams emit protocol-native error events (Chat: \texttt{error} in data; Anthropic: \texttt{event: error}; Responses: \texttt{response.failed}) without any success terminal frame.
